% --- 1. RISK TAXONOMY AND SAFETY GATE ---
% Defines the mapping from tool invocation to safety action.

\begin{equation} \label{eq:risk_level}
\mathcal{R}(T) \in \{ \text{Low, Medium, High, Critical} \}
\end{equation}

\begin{equation} \label{eq:safety_gate}
f_{safe}(T, \Theta) = 
\begin{cases} 
\text{Reject} & \text{if } \exists p \in \mathcal{P}_{reject} : \text{match}(\Theta, p) \\
\text{Suspend} & \text{if } \mathcal{R}(T) \geq \text{High} \\
\text{Execute} & \text{otherwise}
\end{cases}
\end{equation}
where $T$ represents the tool name, $\Theta$ the set of arguments, and $\mathcal{P}_{reject}$ the global list of prohibited execution patterns.

% --- 2. SUPERVISOR ROUTING CONFIDENCE ---
% Formalizes the keyword-confidence selection for specialized agents.

\begin{equation} \label{eq:lexicon_match}
\mathcal{M}(T, \mathcal{L}_A) = \{ w \in \mathcal{L}_A \mid w \text{ is a substring of } T \}
\end{equation}

\begin{equation} \label{eq:routing_confidence}
\mathcal{C}(T, A) = \frac{|\mathcal{M}(T, \mathcal{L}_A)|}{\max_{i \in \mathcal{A}} |\mathcal{L}_i|}
\end{equation}
where $\mathcal{L}_A$ is the domain-specific lexicon for agent $A$ and $\mathcal{A}$ is the set of all available specialized agents.

% --- 3. COMPLEXITY-AWARE ROUTING OBJECTIVE ---
% The optimization objective for selecting the optimal LLM provider tier.

\begin{equation} \label{eq:routing_opt}
p^* = \arg \max_{p \in P} \left( \mathcal{Q}(p, x) - \lambda_C \cdot \mathcal{C}(p, x) - \lambda_L \cdot \mathcal{L}(p, x) \right)
\end{equation}
subject to:
\begin{equation*}
\mathcal{Q}(p, x) \geq \mathcal{Q}_{min}(x)
\end{equation*}
where $\mathcal{Q}$ is reasoning quality, $\mathcal{C}$ is inference cost, $\mathcal{L}$ is latency, and $\lambda$ represents user-defined weighting coefficients.

% --- 4. AGENTIC RAG RERANKING SCORE ---
% Formalizes the two-stage retrieval process.

\begin{equation} \label{eq:rag_sim}
S_{sim}(q, d) = \frac{\mathbf{v}_q \cdot \mathbf{v}_d}{\|\mathbf{v}_q\| \|\mathbf{v}_d\|}
\end{equation}

\begin{equation} \label{eq:rag_final}
S_{final}(q, d) = \Phi_{ce}(q, d) \text{ for } d \in \text{top-}k(\text{arg max } S_{sim})
\end{equation}
where $\Phi_{ce}$ represents the cross-encoder reranking function applied to the top $k$ candidates retrieved via cosine similarity.

% --- 5. STATE TRANSITION FUNCTION ---
% Formalizes the LangGraph workflow state progression.

\begin{equation} \label{eq:state_transition}
\mathcal{S}_{t+1} = \delta(\mathcal{S}_t, \mathcal{A}_t, \mathcal{O}_t)
\end{equation}
where $\mathcal{S}_t$ is the 17-field AgentState at time $t$, $\mathcal{A}_t$ is the action generated by the master orchestrator, and $\mathcal{O}_t$ is the observation retrieved from tool execution.
